\begin{ex}
Đổi: 1 phút = \dots giây.
\shortans{60}
\loigiai{
1 phút = 60 giây.
}
\end{ex}

\begin{ex}
Đổi: 1 thế kỉ = \dots năm.
\shortans{100}
\loigiai{
1 thế kỉ = 100 năm.
}
\end{ex}

\begin{ex}
Đổi: 2 phút = \dots giây.
\shortans{120}
\loigiai{
$2 \times 60 = 120$ giây.
}
\end{ex}

\begin{ex}
Đổi: 5 phút = \dots giây.
\shortans{300}
\loigiai{
$5 \times 60 = 300$ giây.
}
\end{ex}

\begin{ex}
Đổi: 1 phút 30 giây = \dots giây.
\shortans{90}
\loigiai{
$60 + 30 = 90$ giây.
}
\end{ex}

\begin{ex}
Đổi: 2 phút 15 giây = \dots giây.
\shortans{135}
\loigiai{
$120 + 15 = 135$ giây.
}
\end{ex}

\begin{ex}
Đổi: 180 giây = \dots phút.
\shortans{3}
\loigiai{
$180 : 60 = 3$ phút.
}
\end{ex}

\begin{ex}
Đổi: 300 giây = \dots phút.
\shortans{5}
\loigiai{
$300 : 60 = 5$ phút.
}
\end{ex}

\begin{ex}
Đổi: 3 thế kỉ = \dots năm.
\shortans{300}
\loigiai{
$3 \times 100 = 300$ năm.
}
\end{ex}

\begin{ex}
Đổi: 5 thế kỉ = \dots năm.
\shortans{500}
\loigiai{
$5 \times 100 = 500$ năm.
}
\end{ex}

\begin{ex}
Đổi: 10 thế kỉ = \dots năm.
\shortans{1000}
\loigiai{
$10 \times 100 = 1000$ năm.
}
\end{ex}

\begin{ex}
Năm 2000 thuộc thế kỉ thứ bao nhiêu? (Ghi số La Mã hoặc số tự nhiên: thế kỉ 20 ghi 20).
\shortans{20}
\loigiai{
Năm 2000 thuộc thế kỉ XX (20).
}
\end{ex}

\begin{ex}
Năm 2001 thuộc thế kỉ thứ bao nhiêu? (Ghi số tự nhiên).
\shortans{21}
\loigiai{
Năm 2001 thuộc thế kỉ XXI (21).
}
\end{ex}

\begin{ex}
Năm 2024 thuộc thế kỉ thứ bao nhiêu? (Ghi số tự nhiên).
\shortans{21}
\loigiai{
Thế kỉ XXI (21).
}
\end{ex}

\begin{ex}
Năm 1945 thuộc thế kỉ thứ bao nhiêu? (Ghi số tự nhiên).
\shortans{20}
\loigiai{
Năm 1945 thuộc thế kỉ XX (20).
}
\end{ex}

\begin{ex}
Năm 1010 thuộc thế kỉ thứ bao nhiêu? (Ghi số tự nhiên).
\shortans{11}
\loigiai{
Năm 1010 (Lý Thái Tổ dời đô) thuộc thế kỉ XI (11).
}
\end{ex}

\begin{ex}
Năm 938 thuộc thế kỉ thứ bao nhiêu? (Ghi số tự nhiên).
\shortans{10}
\loigiai{
Năm 938 (chiến thắng Bạch Đằng) thuộc thế kỉ X (10).
}
\end{ex}

\begin{ex}
Bác Hồ sinh năm 1890. Năm đó thuộc thế kỉ thứ bao nhiêu? (Ghi số tự nhiên).
\shortans{19}
\loigiai{
Năm 1890 thuộc thế kỉ XIX (19).
}
\end{ex}

\begin{ex}
Thế kỉ thứ 21 bắt đầu từ năm nào?
\shortans{2001}
\loigiai{
Bắt đầu từ năm 2001.
}
\end{ex}

\begin{ex}
Thế kỉ thứ 20 kết thúc vào năm nào?
\shortans{2000}
\loigiai{
Kết thúc vào năm 2000.
}
\end{ex}

\begin{ex}
Một vận động viên chạy 100 m hết 12 giây. Hỏi vận động viên đó chạy 5 lần như thế hết bao nhiêu giây?
\shortans{60}
\loigiai{
$12 \times 5 = 60$ giây.
}
\end{ex}

\begin{ex}
Một nửa phút bằng bao nhiêu giây?
\shortans{30}
\loigiai{
$60 : 2 = 30$ giây.
}
\end{ex}

\begin{ex}
Một phần tư phút bằng bao nhiêu giây?
\shortans{15}
\loigiai{
$60 : 4 = 15$ giây.
}
\end{ex}

\begin{ex}
Một phần năm thế kỉ bằng bao nhiêu năm?
\shortans{20}
\loigiai{
$100 : 5 = 20$ năm.
}
\end{ex}

\begin{ex}
Một phần hai thế kỉ bằng bao nhiêu năm?
\shortans{50}
\loigiai{
$100 : 2 = 50$ năm.
}
\end{ex}
